\section{Experimental setup}

\subsection{Hierarchical deBERTa baseline}

The hierarchical DeBERTa model was implemented using \textbf{Hugging Face transformers} \cite{wolf-etal-2020-transformers} and \textbf{PyTorch}.

\paragraph{Model and Tokenization.}
We used the \texttt{microsoft/mdeberta-v2-base} model as the encoder backbone. The corresponding \texttt{AutoTokenizer} was used with a maximum sequence length of 512 tokens.

\paragraph{Training Configuration.}
The model was trained for a maximum of 10 epochs using the \textbf{Trainer API} with a batch size of 8 and 2 gradient accumulation steps. We used the \textbf{AdamW} optimizer with a learning rate of \(2 \times 10^{-5}\) and a linear warmup schedule for the first 10\% of steps.

\paragraph{Model Selection and Inference.}
Early stopping was configured with a patience of 3 epochs, selecting the best checkpoint based on the \texttt{f1\_sample} score on the evaluation set. An optimal decision threshold was determined via a search on the parent-level predictions of the evaluation set.

\subsection{Zero-Shot Agentic framework}

The following models were utilized:

\begin{itemize}
  \item \textbf{GPT-4o/GPT-4o-mini:} Used as the primary classification agent for narrative and sub-narrative labeling. A temperature of 0 was set to ensure deterministic responses.
  \item \textbf{GPT-4o Mini:} Used as a user proxy agent to relay text to classification agents, chosen for cost efficiency.
  \item \textbf{Zero-Shot/Few-Shot Setup:} Agents classify text based on carefully designed prompts. No fine-tuning was performed.
\end{itemize}

\subsection{Computational Environment}

Experiments were conducted on a machine equipped with: \textbf{Processor:} Core 9 Ultra (22 CPU at 2.5GHz), \textbf{RAM:} 32 GB, \textbf{GPU:} NVIDIA RTX 4070 (8 GB VRAM). However, all LLM inference was performed via API calls to remote servers. The local machine was used primarily to orchestrate API requests, pre-process input text, and handle classification results.

\subsection{Actor-Critic pipeline}

\paragraph{LangGraph implementation.}
For the experiments reported in this paper, the agentic pipeline was implemented using \textbf{LangGraph} \cite{langgraph2024}, a library designed for constructing stateful, multi-actor applications. The reasoning engine for all nodes in the graph (including the Actor and the Critic) was configured to use \textbf{Gemini-2.5-flash}.

\paragraph{Concurrency, batching, and tracing.}
To maximize throughput, we used LangGraph's asynchronous batching capabilities (abatch) and configured it with a maximum concurrency of 20 simultaneous requests. For traceability and prompt development, all LLM interactions were logged using \textbf{LangSmith} \cite{langsmith2024}, producing structured traces of model inputs, outputs, and intermediate reasoning states. These traces were instrumental for debugging and prompt refinement.

